\section{Calibration by Inverse Rendering}
\label{sec:calibration-pipeline}
We calibrate the radiometric camera by inverse rendering: we render
the images of a calibration target with
\cref{eq:rad-flux-with-vignetting} and fit them to the captured RAW
images. The rendered flux depends on the scene as well as on the
camera, so we estimate jointly the poses $\vbTheta_e$, the
projection $\vbTheta_d$, the focus surface $\vbTheta_f$, the aperture
radius $r$ and the vignetting $\vbTheta_v$ of the camera, the
lighting $\vbTheta_l$, the material $\vbTheta_m$ and the shape
$\vbTheta_b$ of the target. We denote all of them by
$\bar{\vbTheta}$. Let $\phi_j(\vbu)$ be the value of the RAW
photosite $\vbu$ in view $j$ and $\sigma_j(\vbu)$ its standard
deviation under the Poisson--Gaussian noise model of the sensor. We
minimize
\begin{equation}
\label{eq:objective}
    \sum_{j}\sum_{\vbu\in\mathcal{S}_j}\!
    \Bigl(\frac{\phi_j(\vbu)-\Phi_j(\vbu;\bar{\vbTheta})}{\sigma_j(\vbu)}\Bigr)^{2}\!
    + \lambda\bigl\lVert \mathrm{D}^3(\hat{\vbd}-\hat{\vbd}^{0})\bigr\rVert^{2}
\end{equation}
over $\bar{\vbTheta}$, where $\mathcal{S}_j$ are photosites drawn
uniformly from those of view $j$ that see the target. The first term
is the negative log-likelihood of the photosites, so that for
$\lambda=0$ the minimizer is the maximum-likelihood estimate of the
camera, the poses and the scene. The second term regularizes the
central spline model only (\cref{sec:optimization}). The remainder of
this section gives what the objective requires: derivatives without
boundary terms, a lighting model, the
reflectance and the shape of the target, and a initialization value.

\subsection{Differentiation without Boundary Terms}
\label{sec:boundary-terms}
\begin{figure}[t]
        \centering
        % fig_boundary.tex — why a boundary that moves with a parameter is missed by fixed samples, and how the calibration removes it
% (hand-written TikZ). Top: the aperture and its radius r (boundary term III). Bottom: an edge of the pattern (I, II).
\usetikzlibrary{arrows.meta}%
\resizebox{\linewidth}{!}{%
\begin{tikzpicture}[x=1cm, y=1cm, >=stealth, font=\footnotesize,
    smp/.style={fill=blue!70!black, draw=none},
    smpn/.style={fill=white, draw=blue!70!black, line width=0.4pt},
    bnd/.style={line width=0.7pt, black},
    moved/.style={line width=0.7pt, red!75!black, densely dashed},
    dim/.style={line width=0.3pt, black!60}]
  % ONE aperture (radius 0.82, growing to 0.98) and ONE pattern of samples inside it, as (distance from the center, angle):
  % the same in the left cell, in the unit disc and in the right cell
  \def\pts{0/0, 0.36/40, 0.36/160, 0.36/280, 0.70/100, 0.70/220, 0.70/340, 0.58/5, 0.64/55, 0.60/185}
  % headers
  \node[anchor=south] at (2.1,1.3) {\textbf{boundary moves}};
  \node[anchor=south] at (6.2,1.3) {\textbf{reformulated}};
  \node[rotate=90, anchor=south] at (0.55,0) {aperture radius $r$};
  \node[rotate=90, anchor=south] at (0.55,-3.0) {pattern edge};
  \draw[dim] (4.05,1.75) -- (4.05,-4.1);
  % ---- top left: the samples are fixed; the ring that the growing rim sweeps holds none
  \fill[red!22, even odd rule] (2.1,0) circle (0.98) (2.1,0) circle (0.82);
  \draw[bnd] (2.1,0) circle (0.82);
  \draw[moved] (2.1,0) circle (0.98);
  \foreach \r/\a in \pts { \fill[smp] ($(2.1,0)+(\a:\r)$) circle (1.4pt); }
  \foreach \r/\a in {1.20/25, 1.18/75, 1.20/130, 1.18/185, 1.20/235, 1.18/285} { \fill[smp] ($(2.1,0)+(\a:\r)$) circle (1.4pt); }
  \node[anchor=west, inner sep=1pt] at (3.18,-0.92) {$\mathcal{A}(r)$};
  \draw[dim] (2.74,-0.52) -- (3.20,-0.88);
  \node[anchor=north] at (2.1,-1.36) {no sample changes};
  % ---- top right: the same aperture and the same samples, placed from the fixed unit disc; they move out with the rim
  \draw[bnd] (4.85,0) circle (0.5);
  \foreach \r/\a in \pts { \fill[smp] ($(4.85,0)+(\a:0.6098*\r)$) circle (1.1pt); }
  \node[anchor=south] at (4.85,0.52) {$\mathcal{D}\ni\vbxi$};
  \draw[->, line width=0.5pt] (5.42,0) -- (5.82,0);
  \node[anchor=north, inner sep=1pt] at (5.4,-0.75) {$\vba=\vbo+r\,\vbxi$};
  \draw[bnd] (7.0,0) circle (0.82);
  \draw[moved] (7.0,0) circle (0.98);
  \foreach \r/\a in \pts {
    \draw[blue!70!black, line width=0.5pt] ($(7.0,0)+(\a:\r)$) -- ($(7.0,0)+(\a:1.1951*\r)$);
    \fill[smp] ($(7.0,0)+(\a:\r)$) circle (1.4pt);
    \draw[smpn] ($(7.0,0)+(\a:1.1951*\r)$) circle (1.4pt);
  }
  \node[anchor=north] at (7.0,-1.36) {every sample moves};
  % ---- bottom left: a sharp edge moves between fixed samples; no sample changes
  \draw[dim, ->] (0.95,-3.75) -- (3.25,-3.75);
  \draw[dim, ->] (0.95,-3.75) -- (0.95,-2.15);
  \node[anchor=east, inner sep=2pt] at (0.95,-2.3) {$\rho$};
  \fill[red!22] (2.0,-3.6) rectangle (2.16,-2.5);
  \draw[bnd] (0.95,-3.6) -- (2.0,-3.6) -- (2.0,-2.5) -- (3.15,-2.5);
  \draw[moved] (2.16,-3.6) -- (2.16,-2.5);
  \foreach \x/\y in {1.2/-3.6, 1.6/-3.6, 2.55/-2.5, 2.95/-2.5} { \fill[smp] (\x,\y) circle (1.5pt); }
  \node[anchor=north] at (2.1,-3.78) {no sample changes};
  % ---- bottom right: the blurred edge; every sample near it changes
  \draw[dim, ->] (4.55,-3.75) -- (7.95,-3.75);
  \draw[dim, ->] (4.55,-3.75) -- (4.55,-2.15);
  \node[anchor=east, inner sep=2pt] at (4.55,-2.3) {$\rho$};
  \draw[bnd, domain=4.55:7.85, samples=60, smooth] plot (\x, {-3.05 + 0.55*tanh((\x-6.05)/0.42)});
  \draw[moved, domain=4.55:7.85, samples=60, smooth] plot (\x, {-3.05 + 0.55*tanh((\x-6.22)/0.42)});
  \foreach \x in {5.45, 5.85, 6.25, 6.65, 7.05} {
    \draw[blue!70!black, line width=0.4pt] (\x, {-3.05 + 0.55*tanh((\x-6.05)/0.42)}) -- (\x, {-3.05 + 0.55*tanh((\x-6.22)/0.42)});
    \fill[smp] (\x, {-3.05 + 0.55*tanh((\x-6.05)/0.42)}) circle (1.5pt);
    \draw[smpn] (\x, {-3.05 + 0.55*tanh((\x-6.22)/0.42)}) circle (1.5pt);
  }
  \node[anchor=north] at (6.25,-3.78) {every sample changes};
\end{tikzpicture}}%

        \caption{Fixed samples (dots) miss a boundary that moves with a parameter (solid black to dashed red): none lies in the region that it sweeps, so the sampled gradient lacks its contribution (left). We remove both boundaries of the calibration (right). \emph{Top:} the samples of the aperture are placed from the fixed unit disc $\mathcal{D}$ and move with its radius $r$. \emph{Bottom:} the blur of the print makes the albedo $\rho$ smooth across an edge of the pattern. Hollow dots: after the change.}
        \label{fig:boundary}
\end{figure}
The derivative of an integral whose domain, or whose integrand's
discontinuities, depend on a parameter has a boundary term, which
automatic differentiation of a sampled estimate misses
(\cref{fig:boundary}, left). Differentiable renderers estimate such
terms by sampling the boundaries
\cite{Li:2018:DMC, Zhang:2020:PSDR, mitsuba-projective}. In the
calibration, three of them arise: The changing aperture radius, border of the calibration target, and sharp texture edges. We avoid these as follows.
%and we remove both by a reformulation (\cref{fig:boundary}, right).
The change of aperture radius is addressed by integrating over the fixed unit disc as described above in \cref{sec:vignetting}. 
%First, the aperture $\mathcal{A}$ depends on its radius $r$. We integrate over the fixed
%unit disc $\mathcal{D}$ (\cref{eq:rad-flux-with-vignetting}): its
%samples $\vbxi$ are placed at $\vba=\vbo+r\vbxi$ and move with the
%radius, so that none of them crosses the rim, and the flux depends
%on $r$ smoothly through the rays of the samples. 
Second, the hard edges of the calibration pattern are discontinuities of the albedo that
move over the aperture and over the pixel with the camera, the pose
and the target shape. We model the print as its design convolved
with a Gaussian, a $10$--$70\,\%$ edge width of $50\,\mathrm{\mu m}$
\cite{rosalen2019cardboards}, so that the albedo is smooth. The
boundary terms of a general scene, from silhouettes, shadows and
the visibility at every bounce of a ray
(\cref{sec:supp-boundaries}), do not arise in the calibration: we only 
render the masked target $\mathcal{M}\subseteq\mathcal{I}$, which is considered smoothly illuminated within a constant scene.

\subsection{Lighting}
\label{sec:lighting}
We represent the light of the environment by spherical harmonics of
degree two \cite{ramamoorthi-2001, peter-pikesloanand2002precomputed},
$L^{\text{SH}}(\vbomega;\vbTheta_l)$, with the coefficients
$\vbTheta_l$ per color channel (\cref{sec:supp-reflectance}). A
calibration target is planar, so every point of it sees the same
hemisphere, and a lighting at infinite distance gives the same
irradiance at every point of the target: any shading across the
target would have to be explained by the vignetting. We therefore
place the lighting on a sphere of finite radius $r_{\text{SH}}$
about the center $\vbb_0$ of the target, the larger of the diagonal
of the target and the largest distance of a camera from $\vbb_0$ at
the start. The radiance that arrives at $\vbx$ from the direction
$\vbomega$ is $L^{\text{SH}}$ in the direction
\begin{equation}
    \hat{\vbomega}(\vbx,\vbomega) \coloneqq \frac{\vbx + \gamma\vbomega - \vbb_0}{\lVert \vbx + \gamma\vbomega - \vbb_0 \rVert}
\end{equation}
from $\vbb_0$ to the point where the ray from $\vbx$ along $\vbomega$
meets the sphere, with $\gamma>0$ such that
$\lVert \vbx + \gamma \vbomega - \vbb_0\rVert = r_{\text{SH}}$. The
lighting then varies over the target and separates from the
vignetting.

\subsection{Reflectance}
\label{sec:reflectance}
Calibration targets usually consist of a pattern printed on background material. To model its appearance properly, we estimate the reflectance distribution of these materials. The radiance that leaves a scene point $\vbx$ in the direction $\vbomega_o$ is \cite{pharr2023physically}
\begin{align}
    &L(\vbx,\vbomega_o)\label{eq:scene-rad}\\
    &\coloneqq \int_{\mathcal{H}^2(\vbx)}f(\vbx, \vbomega_o,\vbomega_i)L(\vbx'(\vbx,\vbomega_i),-\vbomega_i)\cos\theta_i\odif\vbomega_i\nonumber,
\end{align}
where $\mathcal{H}^2(\vbx)$ is the upper hemisphere over $\vbx$,
$\vbx'(\vbx,\vbomega_i)$ the scene or light point visible from $\vbx$
along $\vbomega_i$, $\theta_i$ the angle between $\vbomega_i$ and the
surface normal, and $f$ the bidirectional scattering distribution
function. On the target, the incident radiance is the lighting of
\cref{sec:lighting}. For $f$, we use the Cook--Torrance model
\cite{cookTorrance} with an isotropic GGX microfacet distribution of
roughness $\alpha$ \cite{Heitz2014Microfacet} and the Fresnel term
of a dielectric with the index of refraction $\eta$, and a weight
$\mu\in[0,1]$ between its diffuse part $\rho(\vbx)/\pi$, with the
albedo $\rho$ of the print, and its specular part. We evaluate the
integral by Monte Carlo sampling of the visible microfacet normals
\cite{sphericalCaps} (\cref{sec:supp-reflectance}). Ink fills the
microfacets of the print, so that inked areas are more specular than
blank ones. We therefore blend the roughness and the index of
refraction between a value for ink and one for the blank print by
the albedo, $\vbTheta_m = \{\alpha_{\neg\text{ink}},
\alpha_{\text{ink}},\eta_{\neg\text{ink}}, \eta_{\text{ink}}, \mu\}$.

% Figure of Sec. 5.1 (Detector Bias), declared early so that it floats to the top of the page before the experiments.
\begin{figure*}[t]
        \centering
        % detector-bias-mitsuba.tex — generated by experiments/pgf/detector_bias.py; do not edit.
% The figure: its panels (one PGFPlots picture each, detector-bias-mitsuba/*.tex) stacked in a tabular. The column titles, the row names, the
% renders' heads and the keys' numbers are LaTeX text in the host document's fonts; the keys lie on the panels (one per bias
% panel; Geometric and Radiometric share the one between them). \input it inside a figure*; it needs \usepackage{pgfplots} and
% \pgfplotsset{compat=1.18}. \rcfigdir is the folder this file lies in, as the host document sees it (default fig).
% The six panels share the line width left by the name column (22 pt) and the gaps (33 pt).
\providecommand{\rcfigdir}{fig}%
\begingroup
\ifdefined\rcpanelwd\else\newlength\rcpanelwd\fi\ifdefined\rcpanelht\else\newlength\rcpanelht\fi
\setlength{\rcpanelwd}{\dimexpr(\linewidth-55pt)/6\relax}\setlength{\rcpanelht}{0.66917\rcpanelwd}%
\def\rcpanelwidth{\rcpanelwd}%
\ifdefined\rcflowkeybox\else\newsavebox\rcflowkeybox\fi
\sbox\rcflowkeybox{\scriptsize\input{\rcfigdir/detector-bias-mitsuba/flow-key}}%
\ifdefined\rcflowkeybeyondbox\else\newsavebox\rcflowkeybeyondbox\fi
\sbox\rcflowkeybeyondbox{\scriptsize\input{\rcfigdir/detector-bias-mitsuba/flow-key-beyond}}%
\footnotesize\setlength{\tabcolsep}{0pt}%
\begin{tabular}{@{}c@{\hspace{2pt}}c@{\hspace{10pt}}c@{\hspace{10pt}}c@{\hspace{1.5pt}}c@{\hspace{10pt}}c@{\hspace{1.5pt}}c@{}}
 & Render & \makebox[\rcpanelwd][c]{True Projected vs.\ Detected} & \multicolumn{2}{c@{\hspace{10pt}}}{Calibrated Camera \& Pose} & \multicolumn{2}{c@{}}{Calibrated Camera} \\
 & & & {\scriptsize Geometric} & {\scriptsize Radiometric} & {\scriptsize Geometric} & {\scriptsize Radiometric} \\\noalign{\vskip 2pt}
\raisebox{\dimexpr-0.5\height-0\rcpanelht\relax}[0pt][0pt]{\rotatebox{90}{\footnotesize \begin{tabular}{@{}c@{}}OpenCV +\\AprilGrid\end{tabular}}} & \input{\rcfigdir/detector-bias-mitsuba/opencv-aprilgrid-ocv14-render}\makebox[0pt][r]{\makebox[\rcpanelwd][c]{\raisebox{\dimexpr0.5\rcpanelht-1.5pt\relax}[0pt][0pt]{\tikz[baseline=(k.north)]\node[fill=white,inner sep=1.2pt,font=\scriptsize](k){FOV};}}} & \input{\rcfigdir/detector-bias-mitsuba/opencv-aprilgrid-ocv14-bias}\makebox[0pt][r]{\raisebox{\dimexpr-0.5\rcpanelht+1.5pt\relax}[0pt][0pt]{\tikz[baseline=(k.south)]\node[fill=white,inner sep=1.2pt](k){\usebox\rcflowkeybox\,{\scriptsize $0.24$\,px}};}\kern2pt} & \input{\rcfigdir/detector-bias-mitsuba/opencv-aprilgrid-ocv14-pose-geometric} & \input{\rcfigdir/detector-bias-mitsuba/opencv-aprilgrid-ocv14-pose-radiometric}\makebox[0pt][r]{\makebox[\rcpanelwd][c]{\raisebox{-0.5ex}[0pt][0pt]{\scriptsize\textcolor{black!35}{error not visible}}}}\makebox[0pt][r]{\makebox[\dimexpr2\rcpanelwd+1.5pt\relax][c]{\raisebox{\dimexpr-0.5\rcpanelht+1.5pt\relax}[0pt][0pt]{\tikz[baseline=(k.south)]\node[fill=white,inner sep=1.2pt](k){\usebox\rcflowkeybox\,{\scriptsize $0.39$\,px}};}}}\makebox[0pt][r]{\makebox[\dimexpr2\rcpanelwd+1.5pt\relax][c]{\raisebox{\dimexpr0.5\rcpanelht-1.5pt\relax}[0pt][0pt]{\tikz[baseline=(k.north)]\node[fill=white,inner sep=1.2pt,font=\scriptsize](k){OpenCV-14};}}} & \input{\rcfigdir/detector-bias-mitsuba/opencv-aprilgrid-ocv14-camera-geometric} & \input{\rcfigdir/detector-bias-mitsuba/opencv-aprilgrid-ocv14-camera-radiometric}\makebox[0pt][r]{\makebox[\rcpanelwd][c]{\raisebox{-0.5ex}[0pt][0pt]{\scriptsize\textcolor{black!35}{error not visible}}}}\makebox[0pt][r]{\makebox[\dimexpr2\rcpanelwd+1.5pt\relax][c]{\raisebox{\dimexpr-0.5\rcpanelht+1.5pt\relax}[0pt][0pt]{\tikz[baseline=(k.south)]\node[fill=white,inner sep=1.2pt](k){\usebox\rcflowkeybox\,{\scriptsize $0.54$\,px}};}}}\makebox[0pt][r]{\makebox[\dimexpr2\rcpanelwd+1.5pt\relax][c]{\raisebox{\dimexpr0.5\rcpanelht-1.5pt\relax}[0pt][0pt]{\tikz[baseline=(k.north)]\node[fill=white,inner sep=1.2pt,font=\scriptsize](k){OpenCV-14, global rotation removed};}}} \\\noalign{\vskip 1.5pt}
\raisebox{\dimexpr-0.5\height-0\rcpanelht\relax}[0pt][0pt]{\rotatebox{90}{\footnotesize \begin{tabular}{@{}c@{}}Sch\"ops +\\Siemens stars\end{tabular}}} & \input{\rcfigdir/detector-bias-mitsuba/schops-siemensstars-spline-ds-render}\makebox[0pt][r]{\makebox[\rcpanelwd][c]{\raisebox{\dimexpr0.5\rcpanelht-1.5pt\relax}[0pt][0pt]{\tikz[baseline=(k.north)]\node[fill=white,inner sep=1.2pt,font=\scriptsize](k){Double Sphere};}}} & \input{\rcfigdir/detector-bias-mitsuba/schops-siemensstars-spline-ds-bias}\makebox[0pt][r]{\raisebox{\dimexpr-0.5\rcpanelht+1.5pt\relax}[0pt][0pt]{\tikz[baseline=(k.south)]\node[fill=white,inner sep=1.2pt](k){\usebox\rcflowkeybox\,{\scriptsize $0.067$\,px}};}\kern2pt} & \input{\rcfigdir/detector-bias-mitsuba/schops-siemensstars-spline-ds-pose-geometric} & \input{\rcfigdir/detector-bias-mitsuba/schops-siemensstars-spline-ds-pose-radiometric}\makebox[0pt][r]{\makebox[\dimexpr2\rcpanelwd+1.5pt\relax][c]{\raisebox{\dimexpr-0.5\rcpanelht+1.5pt\relax}[0pt][0pt]{\tikz[baseline=(k.south)]\node[fill=white,inner sep=1.2pt](k){\usebox\rcflowkeybeyondbox\,{\scriptsize $\geq0.053$\,px}};}}}\makebox[0pt][r]{\makebox[\dimexpr2\rcpanelwd+1.5pt\relax][c]{\raisebox{\dimexpr0.5\rcpanelht-1.5pt\relax}[0pt][0pt]{\tikz[baseline=(k.north)]\node[fill=white,inner sep=1.2pt,font=\scriptsize](k){Spline};}}} & \input{\rcfigdir/detector-bias-mitsuba/schops-siemensstars-spline-ds-camera-geometric} & \input{\rcfigdir/detector-bias-mitsuba/schops-siemensstars-spline-ds-camera-radiometric}\makebox[0pt][r]{\makebox[\dimexpr2\rcpanelwd+1.5pt\relax][c]{\raisebox{\dimexpr-0.5\rcpanelht+1.5pt\relax}[0pt][0pt]{\tikz[baseline=(k.south)]\node[fill=white,inner sep=1.2pt](k){\usebox\rcflowkeybeyondbox\,{\scriptsize $\geq0.048$\,px}};}}}\makebox[0pt][r]{\makebox[\dimexpr2\rcpanelwd+1.5pt\relax][c]{\raisebox{\dimexpr0.5\rcpanelht-1.5pt\relax}[0pt][0pt]{\tikz[baseline=(k.north)]\node[fill=white,inner sep=1.2pt,font=\scriptsize](k){Spline, global rotation removed};}}} \\
\end{tabular}
\endgroup

        \caption{The bias of a detector (column 2) is absorbed by a geometrically calibrated camera (columns 3 and 4). With the calibrated poses, the bias is pushed into the camera model (column 3); with the true poses (column 4), it shows that the poses compensate for part of it. Our radiometric refinement does not rely on a detector and removes most of the bias in both evaluations.}
        \label{fig:detector-bias}
\end{figure*}
\subsection{Target Shape}
\label{sec:target-shape}
A printed target is not exactly planar, and its print is not exactly
its design. We chart the board by $\vbb\in\mathcal{B}$, in units of
its width, and model its shape by a height field
$h:\mathcal{B}\rightarrow\mathbb{R}$ along the board normal
$\vbn_{\mathcal{B}}$. In a view with the pose
$(\mathrm{R},\vbt)\in\vbTheta_e$, the chart point $\vbb$ is the scene
point
\begin{equation}
\label{eq:board-surface}
    \vbs(\vbb) \coloneqq \mathrm{R}\bigl(\vbb_{\mathcal{B}} + h(\vbb)\,\vbn_{\mathcal{B}}\bigr) + \vbt,
\end{equation}
where $\vbb_{\mathcal{B}}$ is the point of $\vbb$ on the plane of the
board. The height field is a thin plate spline
\cite{thin-plate-spline} through the heights
$\vbh\in\mathbb{R}^{25}$ at a $5\times5$ grid of nodes. A thin plate
spline is the sum of an affine function and a bending, and affine
heights, an offset and two tilts, describe a plane, which is already
represented by the pose. We therefore restrict the heights to the
pure bendings, $\vbh=\mathrm{B}\vby$, where the columns of
$\mathrm{B}\in\mathbb{R}^{25\times22}$ span the nodal heights whose
spline has no affine part, and optimize $\vbTheta_b\coloneqq\vby$.
The shape then has no gauge freedom with the poses, without a
penalty or a fixed node (\cref{sec:supp-gauge-freedoms}), and only
22 parameters. For the spline camera, we also estimate the in-plane error of the print by
an aspect and a skew, at which the albedo is read.

The scene point of \cref{eq:rad-flux-with-vignetting} is the
intersection of the ray from $\vba$ in the direction
$\vbd\coloneqq\vbpi^{-1}(\vbp,\vba)$ with this surface. We solve for
its chart point: with an orthonormal pair $\vbe_1,\vbe_2$
perpendicular to $\vbd$,
\begin{equation}
\label{eq:board-intersection}
    \vbg(\vbb) \coloneqq
    \begin{pmatrix}\vbe_1^\top\\ \vbe_2^\top\end{pmatrix}
    \bigl(\vbs(\vbb)-\vba\bigr) = \mathbf{0},
\end{equation}
a root of two equations in two unknowns that we find by
Levenberg--Marquardt iterations from the intersection with the plane
of the board. It gives $\vbx(\vbp,\vba)=\vbs(\vbb)$ together with the
chart point $\vbb$ at which the albedo is read. The root is
differentiated \wrt the camera, the pose and the heights by the
implicit function theorem,
$\partial\vbb/\partial\theta=-(\partial\vbg/\partial\vbb)^{-1}\,\partial\vbg/\partial\theta$,
and not through its iterations. The effect of the shape on the
reflectance is small, so the reflectance uses the normal of the
plane.

\subsection{Initialization and Optimization}
\label{sec:optimization}
The projection and the poses start from a geometric calibration
(\cref{sec:supp-starts}), the target as a plane, the focus surface
as the plane parallel to the sensor at the mean distance of the
target, and the aperture radius as $f/(2N)$ with the focal length
$f$ and the f-number $N$ of the lens. The flux is linear in the
vignetting and in the lighting and affine in $\mu$
(\cref{sec:supp-linear}), so they start from one pass of linear
least-squares solves: the lighting with $V\equiv1$, then the
vignetting, then $\mu$. We minimize \cref{eq:objective} with the
Levenberg--Marquardt method \cite{marquardt1963algorithm}, with
$\lambda=0$ for the parametric camera models. The nodes of the central spline at
the border of the lens coverage are hardly constrained by the
photosites. For this model only, the second term regularizes the
change of the nodal directions $\hat{\vbd}\subseteq\bar{\vbTheta}$
from their start $\hat{\vbd}^{0}$ with a P-spline penalty
\cite{eilers1996flexible}, where $\mathrm{D}^3$ takes the third finite
differences, the discrete change of curvature, along both axes of the
grid and $\lambda=0.1$ relative to the data term. It
vanishes for a locally quadratic change, so the penalty leaves a
smooth change of the camera free and suppresses its oscillation, and
the border nodes follow the interior.
No other component is regularized. The parameters stay in their
physical ranges by their parameterization: we optimize $\log r$,
$\log\alpha$, $\log(\eta-1)$ and the logit of $\mu$, the poses on
$SE(3)$ \cite{solà2021microlietheorystate}, and the inverse focus
distances are kept non-negative. The gauge freedoms are removed by
the parameterization (\cref{sec:supp-gauge-freedoms}).
